% rudin_style.pdf: number-first theorem headers ("2.36 Theorem."), as in Rudin or Apostol. No hyperref.
\pdfinfoomitdate=1
\pdftrailerid{}
\pdfsuppressptexinfo=-1
\documentclass[11pt]{article}
\pdfpagewidth=\paperwidth \pdfpageheight=\paperheight % US letter whatever the local TeX default
\usepackage[T1]{fontenc}
\hyphenation{Definition Theorem}
\newcommand{\thmhead}[2]{\par\medskip\noindent\textbf{#1\ \ #2.}\ }
\begin{document}

\noindent\textbf{\Large 2\quad Basic Topology}
\medskip

\thmhead{2.35}{Definition}
\emph{A set is compact if every open cover has a finite subcover.}

\thmhead{2.36}{Theorem}
\emph{If a collection of compact sets has the finite intersection property, its intersection is nonempty.}

\thmhead{2.37}{Corollary}
\emph{A nested sequence of nonempty compact sets has a nonempty intersection.}
\clearpage

\noindent\textbf{\Large 3\quad Numerical Sequences}
\medskip

\thmhead{3.9}{Definition}
\emph{A sequence converges if it is a Cauchy sequence in a complete space.}
The partial sums satisfy
\[ s_n = \sum_{k=1}^{n} a_k \qquad\qquad (3.21) \]
for every index.

By Theorem 2.36 and the remark after it, the claim follows. Using Definition 3.9 we
see that every bounded monotone sequence converges, and the formula (3.21) gives the limit.
\end{document}
